\ssr{Введение}

Цели данной лабораторной работы:
\begin{enumerate}
	\item изучение системных средств порождения процессов и потоков в ОС GNU/LINUX.
\end{enumerate}

\ssr{Выполнение}

\ssection{Задание 2. Программа с двумя дочерними процессами}

Программа, создающая два дочерних процесса с использованием двух вызовов fork(), представлена в листинге~\ref{pr1}. Родительский и оба дочерних процесса должны выводят на свой pid и pid родительского процесса, а также текущее время в формате: часы:минуты:секунды. 

\lstinputlisting[
	style=cpp, 
	caption={Программа, создающая два дочерних процесса},
	label={pr1}
]{program1.c}

\begin{figure}[!htbp]
	\centering
	%\includegraphics[width=0.6\linewidth]{""}
	\caption{Вывод программы}
\end{figure}
\FloatBarrier

\ssection{Задание 3. Программа с двумя дочерними процессами}

Программа, создающая два дочерних процесса с использованием двух вызовов fork(), представлена в листинге~\ref{pr2}. Родительский и оба дочерних процесса должны выводят на свой id и pid своего и родительского процесса и текущее время в формате: часы:минуты:секунды. 

\lstinputlisting[
	style=cpp, 
	caption={Программа, создающая два дочерних процесса},
	label={pr2}
]{program2.c}

\begin{figure}[!htbp]
	\centering
	%\includegraphics[width=0.6\linewidth]{""}
	\caption{Вывод программы}
\end{figure}
\FloatBarrier

\ssection{Задание 4. Модифицированная программа}

Модификация программы из задания 2 таким образом, что второй дочерний процесс после вывода основной информации замещается задачей, создающей два дочерних потока (задание 3) с использованием функции из семейства exec(), представлена в листинге~\ref{pr3}.

\lstinputlisting[
	style=cpp, 
	caption={Модифицированная программа},
	label={pr3}
]{program3.c}

\begin{figure}[!htbp]
	\centering
	%\includegraphics[width=0.6\linewidth]{""}
	\caption{Вывод программы}
\end{figure}
\FloatBarrier

\ssection{Дополнительное задание}

%Получите информацию о системных вызовах, выполняющихся при работе программы, создающей дочерний процесс (системные вызовы требуется отследить как для родительского, так и для дочернего процессов) в двух сценариях:
%1. Родительский и дочерний процесс выводят разные данные (например, различные строковые константы) в стандартный поток вывода;
%2. Родительский процесс, после создания дочернего, выводит отличную от “Hello world!” строку в стандартный поток вывода; дочерний процесс после создания, замещает себя задачей, выводящей на экран строку “Hello world!”.
%Сравнить полученные трассировки системных вызовов и сформулировать выводы по результатам. В отчет включить трассировки для указанных сценариев и выводы по результатам сравнения.

\ssr{Вывод}

В ходе лабораторной работы мы получили 
\begin{itemize}
	\item --
\end{itemize}

\ssr{Контрольные вопросы}

\begin{enumerate}
	\item \textit{}\newline 
	
	\item \textit{В чем различие между процессом и потоком? Какие ресурсы разделяют потоки процесса? Какие ресурсы у каждого потока процесса свои?}\newline 
	
	\item \textit{Что происходит в операционной системе при успешном выполнении fork()?}\newline 
	
	\item \textit{Почему после вызова fork() программа продолжает выполняться с оператора, следующего за вызовом?}\newline 
	
	\item \textit{Какое значение и в какие процессы возвращает функция fork()?}\newline 
	
	\item \textit{Для чего необходимы функции wait() и waitpid()? Что произойдет с дочерним процессом, если родительский процесс завершится до завершения дочернего?}\newline 
	
	\item \textit{Что такое "зомби"-процесс (zombie)? Какими способами можно вывести список всех "зомби"-процессов в системе?}\newline 
	
	\item \textit{Какую задачу решают функции семейства exec? В чем отличие работы exec() от работы fork()?}\newline 
	
	\item \textit{Сохраняется ли PID процесса при замещении его новой программой с помощью exec()?}\newline 
	
	\item \textit{Каков жизненный цикл потока? В каких состояниях может находиться поток?}\newline 
	
	\item \textit{Для чего используется функция pthread\_join()? Что произойдет, если главный поток завершится (например, main вернет управление), а созданные им дочерние потоки еще будут работать?}\newline 
	
\end{enumerate}
